\begin{definition}[Entrywise periodic $Z$-gauge ratio]\label{def:z_gauge_entry}
    \lean{MPSTensor.zGaugeEntry}
    \leanok
    Given complex multiplicity entries $\mu,\nu$, define the entrywise
    $Z$-gauge ratio by
    \[
        z(\mu,\nu) := \mu/\nu.
    \]
\end{definition}

\begin{lemma}[Matched powers imply root-of-unity ratio]\label{lem:z_gauge_entry_pow}
    \lean{MPSTensor.zGaugeEntry_pow_eq_one_of_pow_eq}
    \leanok
    \uses{def:z_gauge_entry}
    If $\mu^m=\nu^m$ and $\nu\neq 0$, then
    \[
        z(\mu,\nu)^m = 1.
    \]
\end{lemma}

\begin{lemma}[Ratio rescales denominator back to numerator]\label{lem:z_gauge_entry_mul}
    \lean{MPSTensor.zGaugeEntry_mul_right}
    \leanok
    \uses{def:z_gauge_entry}
    If $\nu\neq 0$, then
    \[
        z(\mu,\nu) \nu=\mu.
    \]
\end{lemma}

\begin{definition}[Diagonal periodic $Z$-gauge matrix]\label{def:z_gauge_diagonal}
    \lean{MPSTensor.zGaugeDiagonal}
    \leanok
    \uses{def:z_gauge_entry}
    For matched multiplicity entries $(\mu_i)_i$ and $(\nu_i)_i$ on a finite
    index type, define the diagonal matrix
    \[
        Z := \diag\!(z(\mu_i,\nu_i))_i.
    \]
\end{definition}

\begin{lemma}[Diagonal $Z$-gauge has finite order under matched powers]%
    \label{lem:z_gauge_diagonal_pow}
    \lean{MPSTensor.zGaugeDiagonal_pow_eq_one}
    \leanok
    \uses{def:z_gauge_diagonal, lem:z_gauge_entry_pow}
    If $\mu_i^m=\nu_i^m$ and $\nu_i\neq 0$ for all~$i$, then
    \[
        Z^m=\Id.
    \]
\end{lemma}

\begin{lemma}[Diagonal $Z$-gauge transports diagonal multiplicities]\label{lem:z_gauge_diagonal_mul}
    \lean{MPSTensor.zGaugeDiagonal_mul_diagonal}
    \leanok
    \uses{def:z_gauge_diagonal, lem:z_gauge_entry_mul}
    Under $\nu_i\neq 0$ for all~$i$,
    \[
        Z\,\diag(\nu)=\diag(\mu).
    \]
\end{lemma}

\begin{definition}[$Z$-gauge equivalence]
    \label{def:z_gauge_equiv}
    \lean{MPSTensor.ZGaugeEquiv}
    \leanok
    Let $A$ be an $m$-periodic irreducible tensor of bond dimension $D$.
    A global finite-order intertwining relation is given by a matrix
    $Z \in \MN{D}$ satisfying:
    \begin{enumerate}
        \item $Z$ commutes with every tensor matrix: $[A^i, Z] = 0$
              for all physical indices $i$, and
        \item $Z^m = \Id_D$.
    \end{enumerate}
    Two tensors $A, B$ are \emph{$Z$-gauge equivalent} if there exists
    an invertible $Y$ and a $Z$-gauge transformation $Z$ for $A$ such that
    $Z A^i = Y B^i Y^{-1}$ for all~$i$.
    The source equal-case theorem uses the finer block-multiplicity diagonal
    gauges $Z_j$ from Theorem~\ref{thm:periodic_ft}; those multiplicity-space
    matrices are not part of this global relation.
\end{definition}

\begin{remark}[$Z$-gauge is trivial for period-$1$ blocks]
    \label{rem:z_gauge_trivial_period1}
    \lean{MPSTensor.ZGaugeEquiv.of_period_one}
    \leanok
    When $m = 1$ the condition $Z^1 = \Id$ forces $Z = \Id$,
    so $Z$-gauge equivalence reduces to ordinary gauge equivalence
    (Definition~\ref{def:gauge_equiv}).
    The $Z$-gauge is therefore a genuine new degree of freedom
    that appears only for $m > 1$.
\end{remark}

\begin{theorem}[Blocking a $\mathbb Z_m$-gauge gives an ordinary gauge]
    \label{thm:z_gauge_blockTensor_gauge_equiv}
    \lean{MPSTensor.ZGaugeEquiv.blockTensor_gaugeEquiv}
    \leanok
    \uses{def:z_gauge_equiv, def:blocked_tensor, def:gauge_equiv}
    If $A$ and $B$ are $\mathbb Z_m$-gauge equivalent, then their $m$-blocked
    tensors $A^{[m]}$ and $B^{[m]}$ are gauge equivalent in the ordinary sense.
\end{theorem}

\begin{proof}\leanok
    Write the $\mathbb Z_m$-gauge relation as
    $Z A^i = Y B^i Y^{-1}$ with $Z^m = \Id$ and $[A^i, Z] = 0$.
    For a blocked letter, commute one factor of $Z$ past each physical letter in
    the associated length-$m$ word. The total prefactor becomes $Z^m = \Id$,
    so the entire blocked word is related by the same ordinary gauge $Y$.
\end{proof}

\begin{corollary}[Blocking a $\mathbb Z_m$-gauge preserves the MPV family]
    \label{cor:z_gauge_blockTensor_same_mpv}
    \lean{MPSTensor.ZGaugeEquiv.blockTensor_sameMPV}
    \leanok
    \uses{thm:z_gauge_blockTensor_gauge_equiv, thm:gauge_invariance}
    If $A$ and $B$ are $\mathbb Z_m$-gauge equivalent, then the blocked tensors
    $A^{[m]}$ and $B^{[m]}$ have the same MPV family.
\end{corollary}

\begin{proof}\leanok
    Apply Theorem~\ref{thm:z_gauge_blockTensor_gauge_equiv} and then the gauge
    invariance of MPV families.
\end{proof}
